% Generated draft: requires mathematical and source review.
\documentclass[11pt]{article}
\usepackage[margin=25mm]{geometry}
\usepackage{fontspec}
\setmainfont{Latin Modern Roman}
\usepackage{amsmath,amssymb,mathtools,unicode-math}
\setmathfont{Latin Modern Math}
\usepackage{enumitem,needspace,xurl,hyperref}
\setlength{\parindent}{0pt}
\setlength{\parskip}{5pt}
\setlength{\emergencystretch}{3em}
\newcommand{\N}{\mathbb{N}}
\newcommand{\Z}{\mathbb{Z}}
\newcommand{\Q}{\mathbb{Q}}
\newcommand{\R}{\mathbb{R}}
\newcommand{\C}{\mathbb{C}}
\newcommand{\F}{\mathbb{F}}
\newcommand{\A}{\mathbb{A}}
\newcommand{\PP}{\mathbb{P}}
\newcommand{\E}{\mathbb{E}}
\newcommand{\eps}{\varepsilon}
\newcommand{\dd}{\,\mathrm{d}}
\begin{document}
\section*{UnsolvedMath ID: 30003294}
\textbf{Source problem code:} catalogue entry
\subsection{Definitions and mathematical statement}
Under the conventional interpretation of the surface notation, let \(\Sigma_k\) denote a closed, connected, oriented smooth surface of genus \(k\). The supplied background specifies closed fibers; closedness and orientability of the base are understood here from the notation and require confirmation. The record gives no explicit permitted ranges for the genera \(g\) and \(h\).

Let \(\operatorname{Diff}^{+}(\Sigma_g)\) be the group of orientation-preserving smooth diffeomorphisms, and let
\[
\operatorname{Mod}_g=\pi_0\bigl(\operatorname{Diff}^{+}(\Sigma_g)\bigr)
\]
be its mapping class group. Write \(q:\operatorname{Diff}^{+}(\Sigma_g)\to\operatorname{Mod}_g\) for the homomorphism taking a diffeomorphism to its isotopy class. An oriented smooth surface bundle \(p:E\to\Sigma_h\) with fiber \(\Sigma_g\) determines, after choosing a base point and identifying its fiber with \(\Sigma_g\), a monodromy homomorphism
\[
\rho:\pi_1(\Sigma_h)\longrightarrow\operatorname{Mod}_g.
\]
A lift in the category of groups is a group homomorphism \(\widetilde\rho:\pi_1(\Sigma_h)\to\operatorname{Diff}^{+}(\Sigma_g)\) satisfying \(q\circ\widetilde\rho=\rho\). Here the lifting question concerns abstract groups, with no additional continuity requirement.

A smooth bundle is called flat here if it admits local trivializations whose transition diffeomorphisms are locally constant on overlaps, equivalently a flat smooth bundle connection. It is nonflat if no such structure exists on that bundle.

\textbf{Existence question.} Do there exist genera \(g,h\) and an oriented smooth surface bundle
\[
\Sigma_g\longrightarrow E\xrightarrow{\ p\ }\Sigma_h
\]
which is nonflat and whose monodromy satisfies
\[
\nexists\,\widetilde\rho:\pi_1(\Sigma_h)\longrightarrow\operatorname{Diff}^{+}(\Sigma_g)
\quad\text{a group homomorphism with}\quad
q\circ\widetilde\rho=\rho\,?
\]
The intended genus restrictions remain unspecified. Both requirements from the supplied statement are retained; no equivalence between them is asserted here. Nonflatness concerns the bundle itself, without imposing preservation of a distinguished section.
\subsection{Short English statement}
The question asks whether a surface bundle over a surface can have no flat structure and have monodromy that cannot be lifted from mapping classes to actual orientation-preserving diffeomorphisms by a group homomorphism. The intended genus ranges and the conventional interpretation of the base surface require confirmation. Failure to find a flat connection preserving a chosen section does not by itself meet the stated target.
\subsection{Sources}
\begin{itemize}
\item Local UnsolvedMath registry: supplied record with permanent ID \(30003294\), title \emph{Non-Liftable Monodromy of Surface Bundles}, exact source problem code \texttt{catalogue entry}, in the Oberwolfach Reports Open Problems collection. The supplied proposed year is 2016; no proposer is named.
\item Supplied report citation: \emph{Surface Bundles (2017)}, Oberwolfach Reports, DOI 10.4171/owr/2016/56, pp.\ 3149--3195. \url{https://doi.org/10.4171/owr/2016/56}. The supplied literature-source label is ``Oberwolfach Report 56 (2017).'' The citation's year and the year appearing in the DOI are retained as supplied.
\item Supplied literature source: Baykur, \emph{Flat bundles and commutator lengths}. \url{https://arxiv.org/abs/1206.3512}.
\item Literature assessment checked at \texttt{2026-08-21}: this is supplied dated curation, not independently verified. It labels the problem open and reports that Baykur establishes only an approximate statement: certain bundles admit no flat connection making a distinguished section parallel, although the bundles themselves are flat. The supplied assessment says the full closed-fiber lifting problem remains open and requests independent verification of time-sensitive status and cited claims.
\item The registry labels the statement \texttt{corrected\_verified} and reports manual semantic, typographic, and LaTeX review. These are supplied editorial qualifications, not an independent mathematical verification of this draft or of the literature assessment.
\end{itemize}
\end{document}
